\section{Multi-Vertical Use Case}
\label{chap:impl:use-case}

Section~\ref{chap:proposal:use-case} defined the functional and non-functional requirements for a vehicle collision detection and notification network application across the \acs{PPDR} and Automotive vertical sectors. Section~\ref{chap:proposal:capabilities} then identified the testbed capabilities required to support the development and validation of such a use case, and, most importantly, the exposure interfaces through which these capabilities can be made available to third parties. The implementation of the vertical application is presented next, explaining how the capabilities provided by the testbed, already detailed in Sections~\ref{chap:impl:resource_lifecycle} to~\ref{chap:impl:delivery}, are leveraged by the \acs{PoC} to fulfill the application requirements. In doing so, the \acs{PoC} serves as an end-to-end validation of the exposure interfaces and capabilities developed throughout the dissertation, while demonstrating their practical relevance in supporting the implementation and evaluation of future network applications.

Table~\ref{tab:testbed:capabilities} summarizes the testbed's capabilities, highlighting those exercised by this application. Network resources were reserved through the \acs{OSS} by provisioning a dedicated network slice with the uplink/downlink bitrate and latency guarantees needed to sustain continuous video streams from the roadside cameras to the inference server running the collision-detection model. Each camera, a \acs{5G} \acs{UE}, is identified by an \acs{IMSI} and attached to the previous network slice. During operation, network performance is continuously monitored through the testbed's passive monitoring system, presented in Grafana dashboards and also exposed through standardized interfaces, both orchestrated via the \acs{OSS}. The same \acs{OSS} provisions the network exposure components, such as the \acs{NEF} and CAMARA \acsp{API}, integrated with \acs{CAPIF}'s discovery, authentication and authorization workflows, which later support the application's interaction with the \acs{5G} network.

As a \acs{CAPIF} \acs{API} invoker, the application is able to discover and consume the \acsp{API} provided by \acs{NEF} and CAMARA that support the use case. For capabilities supported by both exposure models, the application can use either interface, allowing both alternatives to be exercised. The following sequence diagrams illustrate the interactions involved in the \acs{PoC} implementation for the two alternative \acs{API} exposure paths, prioritizing \acs{NEF}~in~Figure~\ref{chap:impl:use-case:flow:nef} and CAMARA~in~Figure~\ref{chap:impl:use-case:flow:camara}. The experimenter may opt out of the provisioning of the CAMARA component if only the \acs{NEF} is to be utilized, as previously described.

\begin{figure}[htbp]
    \centering
    \includegraphics[width=1\linewidth]{figs/placeholder.png}
    \caption{\acs{PoC} sequence diagram via the \acs{NEF} exposure path}
    \label{chap:impl:use-case:flow:nef}
\end{figure}

The \acs{NEF} interactions are described first (Figure~\ref{chap:impl:use-case:flow:nef}). The network slice is provisioned beforehand through the \acs{OSS}, establishing the baseline connectivity for the application, i.e. the uplink streaming connection between the camera \acs{UE} and the inference server. However, to accommodate the stream's specific \acs{QoS} requirements and their variations over time (e.g. due to network contention), the \acs{PoC} interacts with the \acs{NEF} as part of the application workflow, using the AsSessionWithQoS \acs{API} to establish a \acs{QoS} session and dynamically adjust the connection's \acs{QoS} parameter as requirements change. The application can then continuously monitor the measured uplink throughput, consumed by the live video stream. This is achieved programmatically by creating an analytics subscription to periodically retrieve throughput measurements through the \acs{NEF} AnalyicsExposure \acs{API}, and the Prometheus integration mentioned in Section~\ref{chap:impl:observability}. When the observed metrics breach the \acs{KPI} thresholds for the defined \acs{QoS} level, it requests the enforcement of alternative \acs{QoS} profiles through the same AsSessionWithQoS \acs{API}, keeping the camera streams and detection model working as intended. This adaptation follows a reactive approach, as \acs{QoS} modifications are triggered only after \acs{KPI} violations have been detected. Nonetheless, in the future, \acs{QoS} predictions may be employed for a similar effect. In parallel to all of this, the application tracks pedestrian and vehicle \acs{UE} entering and leaving the monitored area through the \acs{NEF}'s MonitoringEvent \acs{API}.

When a collision is detected, the application triggers two emergency notifications. Pedestrians near the accident are notified via an \acs{SMS} alert delivered through a dedicated mock \acs{SMS} \acs{API}, using their \acs{MSISDN}. In parallel, nearby vehicles are notified through a \acs{DENM} delivered via an \acs{MQTT} broker, using a quadtree-based tiling scheme to relay the message to vehicles by geographic location. The \acs{MQTT} broker was provided separately by the existing 5GAIner infrastructure, whereas the mock \acs{SMS} \acs{API} was specifically developed for this application and included in the \acs{NEF} codebase. To clarify, this choice was made solely for practical reasons, as it avoids the need to deploy an additional component. No such \acs{API} is standardized in any of the \acs{NEF} specifications, something the author is fully aware of.

\begin{figure}[htbp]
    \centering
    \includegraphics[width=1\linewidth]{figs/placeholder.png}
    \caption{\acs{PoC} sequence diagram via the CAMARA exposure path}
    \label{chap:impl:use-case:flow:camara}
\end{figure}

The CAMARA exposure path follows the same overall application workflow, but replaces the specific \acs{NEF} interfaces with the corresponding CAMARA \acsp{API} wherever equivalent capabilities are available. In particular, the \acs{PoC} uses the CAMARA \acs{QoD}, Connectivity Insights Subscriptions, and Device Location \acsp{API} to manage the session's \acs{QoS}, monitor network performance, and track devices in the geofenced area, respectively. Rather than interacting directly with the \acs{NEF}, the network application consumes these higher-level user-friendly CAMARA interfaces, and CAMARA (i.e., the CAMARA translation layer) internally translates the received requests into calls to the corresponding \acs{NEF} AsSessionWithQoS, AnalyticsExposure, and MonitoringEvent \acsp{API}, respectively.

Furthermore, the CAMARA component encapsulates additional logic that would otherwise fall under the responsibility of the network application developer. More specifically, when relying exclusively on the \acs{NEF}, the AnalyticsExposure and MonitoringEvent (\texttt{LOCATION\_REPORTING} event) \acsp{API} periodically retrieve the relevant metrics and location information, respectively, leaving the application to determine whether the measured \acsp{KPI} remain within the defined bounds and the reported locations, within the geofenced area.

With CAMARA, this responsibility shifts to the exposure layer. The application first declares its streaming requirements through the CAMARA Application Profiles \acs{API}, which registers an application profile describing the network conditions needed for the video stream to be processed reliably, namely the minimum upstream and downstream bitrate (i.e. \texttt{targetMinUpstreamRate}, \texttt{targetMinDownstreamRate}). The identifier of this profile is then included in the Connectivity Insights Subscriptions request, so the application no longer has to restate its thresholds whenever it subscribes (unlike AnalyticsExposure). CAMARA uses the referenced profile to evaluate the throughput measurements it receives from the \acs{NEF} and notifies the application only when the defined \acs{KPI} thresholds are breached. A similar event-driven approach applies to geofencing: the application subscribes to location events for the pedestrian and vehicle \acsp{UE} being monitored, and CAMARA notifies it only when a \acs{UE} enters or leaves the defined geofenced area. The application therefore reacts to explicit breach notifications instead of continuously inspecting the raw measurements or location coordinates, which keeps its logic simpler and focused on the use case itself. This is also particularly valuable given the complexity of the \acs{3GPP} data model and its procedures, which is especially evident in the AnalyticsExposure \acs{API}, given the high diversity of analytics event types and payloads.




the rest of the UEs (pedestrians, vehicle UEs are mocked and have their routes preplanned through the capabilities provided by NEFSim

don't forget to mention monitoring event caveat (have to reimplement part of the implementation from the other side)

Other things lack mentioning: 

- mediaMTX

- streaming process (with mediaMTX and ffmpeg)

- roboflow (ML detecting process)


we need like an architecture section to mention the actors and where they are 

The pedestrian and vehicle \acs{UE}s represent bystanders that do not contribute to network traffic and are therefore fully simulated using NEFSim. A population script\footnote{\url{https://github.com/Testbed-TNEI/VehicleDetectionPoC/blob/main/app/scripts/nef_db_init.sh}} is used to register these \acs{UE}s, alongside the actual camera \acs{UE} and the inference server's \acs{IP} address associated with the video stream.

are essentially bystanders that don't have any quote on network traffic, being completely simulated with NEFSim. Mention the 

- print of the dashboards with the stuff of the application

- notification driven 
